\documentclass[11pt, a4paper, reqno]{amsart}

\usepackage{amsmath, amsthm, amssymb, amscd, mathrsfs, mathtools}
\usepackage{graphicx}
\usepackage{hyperref}
\usepackage{booktabs}
\usepackage{color}

% Theorem Environments
\newtheorem{theorem}{Theorem}[section]
\newtheorem{lemma}[theorem]{Lemma}
\newtheorem{proposition}[theorem]{Proposition}
\newtheorem{corollary}[theorem]{Corollary}
\newtheorem{conjecture}[theorem]{Conjecture}
\newtheorem{hypothesis}[theorem]{Hypothesis}

\theoremstyle{definition}
\newtheorem{definition}[theorem]{Definition}
\newtheorem{remark}[theorem]{Remark}
\newtheorem{problem}[theorem]{Open Problem}

% Macros
\newcommand{\R}{\mathbb{R}}
\newcommand{\Q}{\mathbb{Q}}
\newcommand{\Z}{\mathbb{Z}}
\newcommand{\N}{\mathbb{N}}
\newcommand{\K}{\mathbb{Q}((t^{-1}))}
\newcommand{\hHK}{h_{\mathrm{HK}}}
\newcommand{\mutrop}{\mu_{\mathrm{trop}}}
\newcommand{\Vcancel}{V_{\mathrm{cancel}}}

\begin{document}

\title[A Tropical Rigidity Framework for the Erd\H{o}s Conjecture]{A Tropical Rigidity Framework for the Erd\H{o}s Conjecture on Arithmetic Progressions}

\author{kido yoshitaka,gemini}
\address{kanagawa, Japan}
\email{kiddeorn40@gmail.com}
\date{\today}

\begin{abstract}
We propose a novel cross-disciplinary framework aimed at resolving the Erd\H{o}s conjecture on arithmetic progressions. By integrating analytic number theory, the ergodic theory of nilmanifolds, and the tropical geometry of cluster varieties, we translate the arithmetic divergence condition into dynamical rigidity. We formalize the \emph{Leading-Term Non-Vanishing (LTNV) Hypothesis} to establish the exact tropicalization of analytic sequences embedded in formal Laurent series. We then introduce the \emph{Tropical Rigidity Conjecture}, positing that the Host-Kra complexity of a sequence acts as a geometric obstruction (scattering walls) in the tropical fan, suppressing the algebraic entropy $\lambda$ to zero. Conditional on these bridges, the existence of AP-$k$ follows as a topological necessity via Ratner's measure classification and Leibman's polynomial recurrence. Local evidence for $U_3$ (AP-4) and cohomological Maurer-Cartan linking in the wild-type quiver of AP-5 are provided.
\end{abstract}

\maketitle

\section{Introduction and Primary Assumptions}
The Erd\H{o}s conjecture posits that any set $A \subseteq \N$ satisfying $\sum_{a \in A} 1/a = \infty$ contains arithmetic progressions of arbitrary length $k$. We approach this by mapping the sequence into the coordinate ring of a unipotent group $U_k$ endowed with a cluster algebra structure.

To clarify logical dependencies, we operate under the following definitions and assumptions:
\begin{table}[h]
\centering
\begin{tabular}{lp{10cm}}
\toprule
\textbf{Notation/Assumption} & \textbf{Definition \& Application} \\
\midrule
$d_{\log}(A)$ & Logarithmic upper density. Divergence implies $d_{\log}(A) > 0$. \\
$\hHK(\nu)$ & Normalized Gowers $U^k$ norm, measuring Host-Kra complexity. \\
$\lambda$ & Algebraic entropy: $\lambda \coloneqq \limsup_{n \to \infty} \frac{1}{n} \log \deg(\mu^n(X))$. \\
\textbf{Assumption A1} & Subtraction-Free Dynamics \cite{FominZelevinsky2002}. Used in Lemma \ref{lem:ltnv_geom}. \\
\textbf{Assumption A2} & Evaluation in non-Archimedean field $\K$. \\
\textbf{Assumption A3} & Generic initial seeds, establishing a lower bound $c_{\min} > 0$. \\
\bottomrule
\end{tabular}
\end{table}

\section{From Analytic Divergence to Nilfactors}

\begin{lemma}[Logarithmic Measure Construction] \label{lem:measure}
Let $\Omega = \{0, 1\}^\N$. For $A \subseteq \N$ with $\sum_{a \in A} 1/a = \infty$, the weighted averages $\nu_N = \frac{1}{H_N} \sum_{n=1}^N \frac{1}{n} \delta_{T^n 1_A}$ admit a weak* limit $\nu$ by the Banach-Alaoglu theorem. By ergodic decomposition, there exists an ergodic component $\nu_\omega$ satisfying $\nu_\omega([1]_0) \geq d_{\log}(A) > 0$.
\end{lemma}

\begin{lemma}[Divergence implies Non-Trivial $U^k$ Correlation] \label{lem:hk}
Following Host and Kra \cite[Theorem 1.2]{HostKra2005}, a system with vanishing $U^k$ seminorms is weakly mixing. Because $\nu_\omega([1]_0) > 0$, the system avoids pure randomness. The Host-Kra complexity satisfies an auxiliary bound $\hHK(\nu_\omega) \geq f(d_{\log}(A)) > 0$ for some strictly positive function $f$, guaranteeing a non-trivial $k$-step nilfactor.
\end{lemma}

\section{Algebraic Realization and the LTNV Hypothesis}

\begin{hypothesis}[The LTNV Hypothesis] \label{hyp:ltnv}
For a generic arithmetic sequence $A$, the formal expansion of its cluster coordinates $X_i(t) = c_i t^{\alpha_i} + \dots$ satisfies the Leading-Term Non-Vanishing (LTNV) condition under all mutations. Thus, $v(\mu(X)) = \mutrop(v(X))$.
\end{hypothesis}

\begin{lemma}[Geometric Justification of LTNV] \label{lem:ltnv_geom}
Matrix coefficients on unipotent groups $U_k$ possess non-negative integer coefficients. Under Assumption A1, cancellation during mutation $M_1 + M_2$ requires strict algebraic relations, defining a cancellation variety $\Vcancel$. The sequence avoids $\Vcancel$ on a Zariski-open set $\mathcal{U} = \K^m \setminus \Vcancel$, maintaining a global bound $c_{\min} > 0$.
\end{lemma}

\section{The Main Conjecture: Tropical Rigidity}

\begin{conjecture}[Tropical Rigidity Inequality] \label{conj:rigidity}
There exists a constant $C(B, c_{\min}, D) > 0$ such that:
$$\lambda \leq C \cdot (1 - \hHK(\nu_\omega))$$
When $\hHK$ exceeds a critical threshold $h_0$, the density of scattering walls strictly enforces $\lambda = 0$.
\end{conjecture}

\section{Conditional Resolution of the Erd\H{o}s Conjecture}

\begin{theorem}[Conditional Existence of AP-$k$] \label{thm:main}
Assuming Hypothesis \ref{hyp:ltnv} and Conjecture \ref{conj:rigidity}, any sequence $A \subseteq \N$ with $\sum 1/a = \infty$ contains AP-$k$.
\end{theorem}
\begin{proof}
By Lemma \ref{lem:hk}, divergence forces $\hHK > 0$. By Conjecture \ref{conj:rigidity}, this structure suppresses algebraic entropy to $\lambda = 0$, ensuring the orbit $u(n)\Gamma$ is generated by a strictly unipotent flow. By Ratner's Theorem \cite[Theorem 1]{Ratner1991}, the invariant measure is the Haar measure $\mu_H$ on an algebraic sub-nilmanifold $H$. By Leibman's polynomial recurrence theorem \cite{Leibman2005}, the orbit intersects its shifts, yielding AP-$k$.
\end{proof}

\section{Local Evidence: Wild-Type Scattering Cohomology (AP-5)}
For $k=5$, the cluster algebra is of wild type. However, the $U^4$-nilfactor defines a Hamiltonian invariant $H = (x_1^2 + x_3^2 + x_2 x_4 + 1)/(x_1 x_2 x_3 x_4)$. Following Gross-Siebert \cite[Theorem 1.4]{GrossSiebert2011}, this defines a 1-cohomology class $[\omega] \in H^1(\mathcal{D}, \Z)$ ensuring the path-ordered product satisfies the Maurer-Cartan condition ($\prod_L \theta_L = \mathrm{Id}$). This cohomologically links the infinite dense walls into a periodic invariant cylinder, enforcing $\lambda = 0$.

\section{Open Problems}
\begin{problem}
Compute the defining polynomial equations of $\Vcancel$ for the $U_4$ nilmanifold and prove that the Laurent embeddings of positive density sequences admit infinitesimal perturbations to remain strictly within $\mathcal{U}$.
\end{problem}
\begin{problem}
Prove analytically that the 1-cohomology class defined by the level sets of $H_{\mathrm{AP5}}$ satisfies the global Maurer-Cartan condition across the entire Gross-Siebert scattering diagram.
\end{problem}

\bibliographystyle{alpha}
\bibliography{references}

\end{document}